\section{把结果联系起来}
\label{sec:discussion}
\subsection{数学结果的含义}
\label{subsec:implications}
测量结果说明了从一项检查中恢复信息的限度。正确程序与错误程序若得到相同的记录描述，只使用这份描述的规则就必须同样处理两者。在八个可执行布尔样例中，把两道检查通过解释为接受会错判五个；仍然只用这两道检查的最优强制二分类规则，也至少错判三个。要区分这些程序，需要增加观察，而不只是重新计算已有输出。

过程结果说明，概率模型必须同时交代观察规则。条件通过概率恒定，与后续步骤剩余路径较少并不矛盾。首次通过后停止、由于其他原因退出，以及只保留成功路径，会产生不同的后果。这些结论来自指定模型；重建出一条历史路径，还不足以确认模型假设。

识别结果给出了判断是否还需要证据的准则：如果每个允许的补全都给出相同差值，现有信息便确定了它；否则，增加观察才可能缩小范围。夏普利分摊能够概括完整比较，但不会带来新的观察。

\subsection{为什么两种历史趋势可以同时成立}
\label{subsec:two-trends}
在选定的候选修改边上，以任务等权计算的记录通过变化为正。在固定失败起点路径上，后续被观察步骤的通过比例却低于首步。两句话使用不同的对象与筛选条件，没有任何代数规则要求它们方向一致。

例如，设十个未解决任务中，八个在第一步通过，余下两个里有一个在第二步通过。通过比例从 $8/10=80\%$ 降为 $1/2=50\%$，已有通过记录的任务却从八个增加到九个。第二组恰好由前一步没有通过的任务组成，因此其比例较低，并不能说明第二次修改伤害了任务。这是一个说明性例子。

真实项目还有更多复杂性：同一个任务可能有多条路径，候选和审核条件可能改变，路径也可能离开规定的图结构。因此，实际的任务配对分析与过程分析分别保持自己的定义，不把说明性例子当作已经拟合的数据模型。它们得到的正、负差值都涉及记录结果，而不是独立确认的数学质量。

\subsection{逻辑、统计与题意上的不确定性}
\label{subsec:uncertainty}
报告使用了三种不同的不确定性。把它们分开，公式的含义就更容易理解。
\begin{itemize}[leftmargin=*]
\item \textbf{尚未观察的比较：}一个分数缺失，所以多张表都可能与现有信息相容。在有效约束下补全这些表，可以得到答案的逻辑范围。
\item \textbf{统计参考计算：}把保留任务视为来自某种明确、可以设想重复的抽样机制。标准误或重抽样区间描述在该机制下的变化；能否适用于观察任务之外，取决于相关假设。
\item \textbf{尚未解决的题意解释：}我们没有独立参考来判断候选是否表达了预期数学。针对通过标签得到再窄的抽样区间，也不能补上这个参考。
\end{itemize}
布尔对照已列举全部输入，并观察到相应输出；其局限在于，这个任务上的结论能够推广到多大范围。因此，有限样例的错误下界、缺失格的逻辑界和抽样区间具有不同含义。

\subsection{后续比较研究的条件}
\label{subsec:future-design}
如果以后要比较题意评价能力，先确定任务怎样选取，以及具体测量哪个结果。例如，可以请独立裁定者按明确参考判断候选是否满足某一条要求。向被比较的评价器提供同一候选、同一要求和同一范围的证据；在任务内部比较它们的答案，保留无法判断的情况，并将同任务的重复观察放在一起。构造反例集仍然有用，但它与独立抽样的评估集回答不同问题。

还可以用一个简单规划关系，理解精度与独立任务数量的联系。设 $h>0$ 是任务配对平均差的近似 95\% 区间所期望的半宽。设 $\sigma_D$ 是拟议抽样总体中任务差值的标准差，应来自合适的先导研究，或采用有依据的保守上界。设 $S_{\rm plan}$ 是计划纳入的任务数量，不是运行次数或修改边数。在任务独立抽样的条件下，标准误近似为 $\sigma_D/\sqrt{S_{\rm plan}}$。正态参考下的 95\% 区间，其半宽近似为标准误的 1.96 倍。

要求半宽不大于 $h$，依次得到
\[
1.96\,\frac{\sigma_D}{\sqrt{S_{\rm plan}}}\le h,
\qquad
\sqrt{S_{\rm plan}}\ge\frac{1.96\,\sigma_D}{h},
\qquad
S_{\rm plan}\ge\frac{1.96^2\sigma_D^2}{h^2}.
\]
例如，若仅为说明而取标准差 0.20、目标半宽 0.05，右侧约为 61.47，在这个近似下提示至少需要 62 个任务。这些是规划时设定的输入，不是本研究给出的估计值；62 也不是已经对本项目论证充分的建议数量。小样本修正、依赖关系、参考质量和实际抽样设计都可能改变计算。在这个简单关系中，把目标半宽减半，大约需要四倍的独立任务；重复运行同一候选不会创造这些任务。

后续过程研究还需要事先确定：路径何时开始、哪些下一步符合观察条件、观察何时停止，以及退出怎样记录。这些决定应当使用下一次结果出现之前的信息。随后，可在隔离的任务组或较晚时间段检验预测，回答一个明确的实际使用问题。本报告提供可复算的历史描述和具体测量核查，帮助设计这种研究，但不声称后续研究已经完成。
